% !TeX program = lualatex
\documentclass[12pt,aspectratio=169]{beamer}
\usepackage[utf8]{inputenc}
\usepackage[T1]{fontenc}
\usepackage{amsmath,amsfonts,amssymb}
\usepackage{graphicx}
\usepackage{xcolor}
\usepackage{hyperref}
\usepackage{anyfontsize}
\usepackage{lmodern}

% ====== EXTREME FONT REDUCTION ======
\renewcommand{\tiny}{\fontsize{4}{5}\selectfont}
\renewcommand{\scriptsize}{\fontsize{5}{6}\selectfont}
\renewcommand{\footnotesize}{\fontsize{6}{7}\selectfont}
\renewcommand{\small}{\fontsize{7}{8}\selectfont}
\renewcommand{\normalsize}{\fontsize{8}{9}\selectfont}
\renewcommand{\large}{\fontsize{9}{10}\selectfont}
\renewcommand{\Large}{\fontsize{10}{11}\selectfont}
\renewcommand{\LARGE}{\fontsize{11}{13}\selectfont}
\renewcommand{\huge}{\fontsize{12}{14}\selectfont}
\renewcommand{\Huge}{\fontsize{14}{17}\selectfont}

% ====== COLOR THEME ======
\definecolor{redcorrect}{RGB}{200,30,30}
\definecolor{bluecorrect}{RGB}{30,80,200}
\definecolor{greencheck}{RGB}{0,150,50}
\definecolor{lightgray}{RGB}{245,245,245}
\definecolor{darkgray}{RGB}{50,50,50}
\definecolor{softyellow}{RGB}{255,248,200}
\definecolor{steppurple}{RGB}{150,50,200}
\definecolor{deepblue}{RGB}{20,60,150}
\definecolor{darkred}{RGB}{150,20,20}

\setbeamercolor{title}{fg=white,bg=redcorrect}
\setbeamercolor{frametitle}{fg=white,bg=redcorrect}
\setbeamercolor{block title}{fg=white,bg=bluecorrect}
\setbeamercolor{block body}{fg=darkgray,bg=lightgray}
\setbeamercolor{normal text}{fg=darkgray}
\usecolortheme{default}

% ====== CUSTOM COMMANDS ======
\newcommand{\correct}[1]{\textcolor{greencheck}{\textbf{\checkmark}} #1}
\newcommand{\keypoint}[1]{\textcolor{deepblue}{\textbf{#1}}}
\newcommand{\hardconcept}[1]{\textcolor{darkred}{\textbf{#1}}}
\newcommand{\tipbox}[1]{\fcolorbox{bluecorrect}{softyellow}{\parbox{0.88\textwidth}{\centering #1}}}
\newcommand{\stepbox}[1]{%
	\begin{center}
		\fcolorbox{darkgray}{white}{%
			\parbox{0.90\textwidth}{\vspace{0.05cm} #1 \vspace{0.05cm}}%
		}
	\end{center}
}

\begin{document}
	
	% ================================================================
	% SLIDE 1: TITLE
	% ================================================================
	\begin{frame}
		\centering
		\vspace{0.8cm}
		{\fontsize{16}{19}\selectfont \textbf{AML / CFT Test Review}} \\[0.15cm]
		{\fontsize{12}{14}\selectfont Advanced Concepts \& Detailed Analysis}\\[0.3cm]
		{\fontsize{9}{11}\selectfont Understanding the \textit{Regulatory Principles} Behind Each Answer}\\[0.3cm]
		{\fontsize{8}{10}\selectfont Compliance Training Department}
		\vspace{0.2cm}
		\rule{4cm}{0.5pt}
	\end{frame}
	
	% ================================================================
	% SLIDE 2: Screenshot 1 — AML/CFT Program & Controls (Q4)
	% ================================================================
	\begin{frame}
		\frametitle{\fontsize{11}{13}\selectfont 1. Corporate Onboarding \& EDD}
		\begin{block}{\fontsize{8}{10}\selectfont Subject: Design Your AML/CFT Program and Controls}
			\scriptsize
			A financial institution is onboarding a new corporate client. During review, the onboarding team identifies that the company's ownership structure includes multiple layers across different jurisdictions, and one of the shareholders is a nominee company based in a recognized secrecy haven. What should be the onboarding team's next step?
		\end{block}
		\vspace{0.05cm}
		\stepbox{\scriptsize
			\keypoint{Correct Answer:} Escalate the case for enhanced due diligence and verify the identity of the ultimate beneficial owner.
		}
		\vspace{0.05cm}
		\scriptsize
		\hardconcept{Core Concept:} This is a direct application of \textbf{FATF Recommendation 10} (Customer Due Diligence) and \textbf{Recommendation 24} (Transparency of Legal Persons). When a client presents multiple ownership layers and nominee structures, the risk of money laundering increases significantly. The Financial Action Task Force requires that financial institutions identify the natural person who ultimately owns or controls the client.
		\vspace{0.05cm}
		\tipbox{\scriptsize \textbf{Real-World Application:} Always escalate complex structures with secrecy elements. The UBO must be identified before onboarding. Failure to do so can result in regulatory fines under \textbf{AMLD6} and \textbf{US Patriot Act Section 311} designations.}
	\end{frame}
	
	% ================================================================
	% SLIDE 3: Screenshot 2 — AML/CFT Program & Controls (Q5)
	% ================================================================
	\begin{frame}
		\frametitle{\fontsize{11}{13}\selectfont 2. AFC Regulatory Reporting}
		\begin{block}{\fontsize{8}{10}\selectfont Subject: Design Your AML/CFT Program and Controls}
			\scriptsize
			Which of the following practices best support accurate and consistent AFC regulatory reporting? (Select Two.)
		\end{block}
		\vspace{0.05cm}
		\stepbox{\scriptsize
			\keypoint{Correct Answers:} 
			\par Implementing a centralized transaction monitoring system.
			\par Regular data cleansing of KYC and payment data sources.
		}
		\vspace{0.05cm}
		\scriptsize
		\hardconcept{Core Concept:} Regulatory reporting accuracy depends on \textbf{data integrity}. A centralized system ensures a single source of truth, eliminating silos that cause inconsistent reporting. Data cleansing addresses the \textbf{GIGO (Garbage In, Garbage Out)} principle. Under \textbf{FATF Recommendation 15} (New Technologies), firms must ensure their systems produce reliable data for suspicious transaction reporting (STR) and currency transaction reporting (CTR).
		\vspace{0.05cm}
		\tipbox{\scriptsize \textbf{Key Insight:} The \textbf{MIFID II} and \textbf{EMIR} frameworks also require accurate trade reporting. The same data quality principles apply across all regulatory reporting obligations.}
	\end{frame}
	
	% ================================================================
	% SLIDE 4: Screenshot 3 — AML Risk Management (Q1)
	% ================================================================
	\begin{frame}
		\frametitle{\fontsize{11}{13}\selectfont 3. Essential Training Components}
		\begin{block}{\fontsize{8}{10}\selectfont Subject: AML Risk Management Program Components and Duties}
			\scriptsize
			Which of the following are essential components of an ongoing compliance training program? (Select Three.)
		\end{block}
		\vspace{0.05cm}
		\stepbox{\scriptsize
			\keypoint{Correct Answers:}
			\par Customizing training content based on employee roles.
			\par Incorporating updates on recent regulatory changes.
			\par Including metrics to assess participant comprehension.
		}
		\vspace{0.05cm}
		\scriptsize
		\hardconcept{Core Concept:} \textbf{FATF Immediate Outcome 4} requires that financial institutions have competent staff who understand AML/CFT obligations. Role-based training aligns with the \textbf{risk-based approach} — frontline staff need different knowledge than investigators. Regulatory updates address the evolving nature of money laundering typologies. Comprehension metrics satisfy \textbf{Regulatory expectation for training effectiveness} under frameworks like the \textbf{UK Financial Conduct Authority's SYSC 6.1.4} and \textbf{US Bank Secrecy Act} requirements.
		\vspace{0.05cm}
		\tipbox{\scriptsize \textbf{Best Practice:} Implement a training management system (LMS) that tracks completion and comprehension scores. This provides audit evidence of training effectiveness.}
	\end{frame}
	
	% ================================================================
	% SLIDE 5: Screenshot 4 — AML Risk Management (Q2)
	% ================================================================
	\begin{frame}
		\frametitle{\fontsize{11}{13}\selectfont 4. Second Line AFC Officer}
		\begin{block}{\fontsize{8}{10}\selectfont Subject: AML Risk Management Program Components and Duties}
			\scriptsize
			What is a responsibility of a second line Anti-Financial Crime (AFC) officer when developing new controls to detect money laundering?
		\end{block}
		\vspace{0.05cm}
		\stepbox{\scriptsize
			\keypoint{Correct Answer:} Balancing regulatory requirements with data protection obligations.
		}
		\vspace{0.05cm}
		\scriptsize
		\hardconcept{Core Concept:} The \textbf{Three Lines of Defense} model positions the second line as the compliance function that oversees and challenges the first line. The balance between AML and data protection is governed by \textbf{GDPR Article 6} (lawfulness of processing), \textbf{GDPR Article 9} (processing of special categories of data), and \textbf{GDPR Article 5(1)(b)} (purpose limitation). Under \textbf{EU AMLD5}, there is an explicit requirement to balance these obligations. The second line must ensure that data collected for AML purposes is not overused or retained longer than necessary, while still being sufficient for effective monitoring.
		\vspace{0.05cm}
		\tipbox{\scriptsize \textbf{Key Principle:} The second line must document the \textbf{legal basis} for data processing and conduct \textbf{Data Protection Impact Assessments (DPIAs)} for high-risk AML systems.}
	\end{frame}
	
	% ================================================================
	% SLIDE 6: Screenshot 5 — AML Risk Management (Q3)
	% ================================================================
	\begin{frame}
		\frametitle{\fontsize{11}{13}\selectfont 5. Training Program Challenges}
		\begin{block}{\fontsize{8}{10}\selectfont Subject: AML Risk Management Program Components and Duties}
			\scriptsize
			Which of the following challenges can impede the effectiveness of compliance training programs? (Select Three.)
		\end{block}
		\vspace{0.05cm}
		\stepbox{\scriptsize
			\keypoint{Correct Answers:}
			\par Poor alignment between compliance risks and controls explained in training.
			\par Lack of customization for different roles.
			\par Outdated content that does not reflect current threats.
		}
		\vspace{0.05cm}
		\scriptsize
		\hardconcept{Core Concept:} This relates to the \textbf{adult learning theory} applied to compliance. Misalignment between risks and controls creates a \textbf{knowledge-action gap} — staff know the theory but don't know how to apply it. \textbf{Kirkpatrick's Four Levels of Training Evaluation} suggests that without relevance (Level 1) and application (Level 3), training is ineffective. Regulatory bodies like the \textbf{US Financial Crimes Enforcement Network (FinCEN)} and \textbf{UK National Crime Agency (NCA)} emphasize that training must be \textbf{risk-based} and \textbf{role-specific}. Outdated content creates \textbf{compliance drift} — staff learn obsolete procedures.
		\vspace{0.05cm}
		\tipbox{\scriptsize \textbf{Recommended Action:} Conduct a \textbf{Training Needs Analysis (TNA)} annually and update content based on regulatory changes and internal audit findings.}
	\end{frame}
	
	% ================================================================
	% SLIDE 7: Screenshot 6 — AML Risk Management (Q1, second set)
	% ================================================================
	\begin{frame}
		\frametitle{\fontsize{11}{13}\selectfont 6. Business Support Role}
		\begin{block}{\fontsize{8}{10}\selectfont Subject: AML Risk Management Program Components and Duties}
			\scriptsize
			If business development engages new clients without applying risk-based onboarding procedures, what role should business support (within the first line) take?
		\end{block}
		\vspace{0.05cm}
		\stepbox{\scriptsize
			\keypoint{Correct Answer:} Pause client access and ensure procedures are back in scope before processing.
		}
		\vspace{0.05cm}
		\scriptsize
		\hardconcept{Core Concept:} This tests the \textbf{operationalization of the three lines of defense}. The first line (business support) must enforce \textbf{controls within the process}. Pausing access is an example of a \textbf{preventative control} that stops non-compliant activities before they create risk. Under \textbf{FATF Recommendation 10}, a financial institution must not establish a business relationship until CDD measures are complete. This is also consistent with \textbf{the Principle of Proportionality} — the response is proportionate to the risk of allowing an unvetted client access.
		\vspace{0.05cm}
		\tipbox{\scriptsize \textbf{Real-World:} Business support acts as the \textbf{gatekeeper}. If business development bypasses procedures, the escalation should go to the \textbf{MLRO (Money Laundering Reporting Officer)}.}
	\end{frame}
	
	% ================================================================
	% SLIDE 8: Screenshot 7 — AML Risk Management (Q2, second set)
	% ================================================================
	\begin{frame}
		\frametitle{\fontsize{11}{13}\selectfont 7. GDPR \& AML Investigations}
		\begin{block}{\fontsize{8}{10}\selectfont Subject: AML Risk Management Program Components and Duties}
			\scriptsize
			To maintain GDPR compliance while conducting AML investigations, what must the second line do when suspicious activity involves customer data?
		\end{block}
		\vspace{0.05cm}
		\stepbox{\scriptsize
			\keypoint{Correct Answer:} Collaborate with Legal to evaluate data use case and retention legality.
		}
		\vspace{0.05cm}
		\scriptsize
		\hardconcept{Core Concept:} This addresses the \textbf{intersection of data protection and AML/CFT obligations}. Under \textbf{GDPR Article 6(1)(c)} and \textbf{Article 6(1)(e)}, processing for compliance with a legal obligation is lawful. However, under \textbf{GDPR Article 5(1)(b)} (purpose limitation) and \textbf{GDPR Article 25} (data protection by design), the institution must ensure that data is not used beyond what is necessary. The \textbf{EDPB Guidelines on AML/CFT} emphasize that AML data processing must be \textbf{proportionate}. Legal collaboration ensures that the institution can demonstrate \textbf{accountability} under \textbf{GDPR Article 5(2)}.
		\vspace{0.05cm}
		\tipbox{\scriptsize \textbf{Key Principle:} Always document the \textbf{legal basis} for processing and the \textbf{retention period}. This is essential for the institution's \textbf{Record of Processing Activities (ROPA)}.}
	\end{frame}
	
	% ================================================================
	% SLIDE 9: Screenshot 8 — Technology for Monitoring (Q4)
	% ================================================================
	\begin{frame}
		\frametitle{\fontsize{11}{13}\selectfont 8. Change Management Protocol}
		\begin{block}{\fontsize{8}{10}\selectfont Subject: Technology for Ongoing Monitoring and Investigations}
			\scriptsize
			The internal audit team reports that no documentation exists for changes made to transaction monitoring scenarios in the past six months. What should compliance leadership do first?
		\end{block}
		\vspace{0.05cm}
		\stepbox{\scriptsize
			\keypoint{Correct Answer:} Formalize a change management protocol and require documentation moving forward.
		}
		\vspace{0.05cm}
		\scriptsize
		\hardconcept{Core Concept:} This is about \textbf{IT Governance} and \textbf{change management control}. Under \textbf{FATF Recommendation 15} (New Technologies) and \textbf{US FFIEC BSA/AML Examination Manual}, institutions must maintain effective technology controls. Undocumented changes create \textbf{model risk} — the risk that the monitoring system is not performing as intended. This is particularly critical under \textbf{SR 11-7} (Supervisory Guidance on Model Risk Management). The change management protocol should include \textbf{approval workflows}, \textbf{testing requirements}, and \textbf{audit trails}.
		\vspace{0.05cm}
		\tipbox{\scriptsize \textbf{Best Practice:} Implement a \textbf{change management system} (e.g., ITIL) with version control. All changes must be approved by the \textbf{MLRO} and \textbf{Head of Compliance}.}
	\end{frame}
	
	% ================================================================
	% SLIDE 10: Screenshot 9 — Technology for Monitoring (Q5)
	% ================================================================
	\begin{frame}
		\frametitle{\fontsize{11}{13}\selectfont 9. Blockchain Tracing Challenges}
		\begin{block}{\fontsize{8}{10}\selectfont Subject: Technology for Ongoing Monitoring and Investigations}
			\scriptsize
			Which of the following are challenges in tracing illicit activity on blockchain networks? (Select Three.)
		\end{block}
		\vspace{0.05cm}
		\stepbox{\scriptsize
			\keypoint{Correct Answers:}
			\par Use of privacy coins such as Monero and Zcash.
			\par Mixing services that break links between transactions.
			\par Rapid fund movement across multiple addresses.
		}
		\vspace{0.05cm}
		\scriptsize
		\hardconcept{Core Concept:} These challenges relate to \textbf{anonymity-enhancing technologies} and \textbf{transaction layering}. Privacy coins use \textbf{ring signatures} (Monero) or \textbf{zero-knowledge proofs} (Zcash) to obscure the sender, receiver, and amount. Mixing services (e.g., Tornado Cash) use \textbf{peer-to-peer obfuscation} to break the chain of custody. Rapid movement across addresses is a \textbf{layering technique} — moving funds quickly across multiple wallets to confuse blockchain analytics tools. These are recognized by \textbf{FATF Recommendation 16} (Wire Transfer Rules) and the \textbf{FATF Report on Virtual Assets} as key challenges.
		\vspace{0.05cm}
		\tipbox{\scriptsize \textbf{Real-World Insight:} Blockchain analytics tools (Chainalysis, Elliptic) use \textbf{heuristics} to try to overcome these challenges, but they remain significant obstacles.}
	\end{frame}
	
	% ================================================================
	% SLIDE 11: Screenshot 10 — Technology for Monitoring (Q1)
	% ================================================================
	\begin{frame}
		\frametitle{\fontsize{11}{13}\selectfont 10. Payment Screening System}
		\begin{block}{\fontsize{8}{10}\selectfont Subject: Technology for Ongoing Monitoring and Investigations}
			\scriptsize
			A global bank with a high transaction volume is updating its payment screening system. Its goal is to minimize false positives while maintaining regulatory compliance. Which configuration of tools would best achieve this?
		\end{block}
		\vspace{0.05cm}
		\stepbox{\scriptsize
			\keypoint{Correct Answer:} Use a hybrid system combining fuzzy logic, behavioral AI models, and rule-based filters to scrutinize real-time payments.
		}
		\vspace{0.05cm}
		\scriptsize
		\hardconcept{Core Concept:} This addresses \textbf{false positive reduction strategies}. False positives are a major operational cost and can cause \textbf{desensitization} to true alerts. The hybrid approach leverages:
		\par \textbf{Fuzzy Logic}: Handles name variations and typographical errors (e.g., "Mohamed" vs "Muhammad").
		\par \textbf{Behavioral AI}: Learns normal transaction patterns and flags deviations (unsupervised learning).
		\par \textbf{Rule-Based Filters}: Enforces regulatory requirements (e.g., OFAC SDN list matching).
		This is consistent with \textbf{ISO 20022} and \textbf{SWIFT Payment Controls} guidelines. It also aligns with \textbf{FATF Recommendation 14} on real-time monitoring.
		\vspace{0.05cm}
		\tipbox{\scriptsize \textbf{Key Takeaway:} A \textbf{layered defense} reduces false positives while maintaining detection effectiveness.}
	\end{frame}
	
	% ================================================================
	% SLIDE 12: Screenshot 11 — Technology for Monitoring (Q2)
	% ================================================================
	\begin{frame}
		\frametitle{\fontsize{11}{13}\selectfont 11. Blockchain Tracing Technologies}
		\begin{block}{\fontsize{8}{10}\selectfont Subject: Technology for Ongoing Monitoring and Investigations}
			\scriptsize
			Which of the following technologies or methods support blockchain transaction tracing activities? (Select Two.)
		\end{block}
		\vspace{0.05cm}
		\stepbox{\scriptsize
			\keypoint{Correct Answers:}
			\par Wallet clustering algorithms.
			\par Chain visualization dashboards.
		}
		\vspace{0.05cm}
		\scriptsize
		\hardconcept{Core Concept:} \textbf{Wallet clustering} is a \textbf{graph-based algorithm} that groups cryptocurrency addresses that are likely controlled by the same entity. This uses \textbf{common spending heuristics} (addresses that are inputs to the same transaction) and \textbf{co-spending patterns}. \textbf{Chain visualization} creates a \textbf{graph database} of transactions that investigators can navigate to see flows of funds. These methods are used by \textbf{Chainalysis}, \textbf{Elliptic}, and \textbf{CipherTrace}. The underlying principle is \textbf{network analysis} — treating the blockchain as a graph of addresses and transactions.
		\vspace{0.05cm}
		\tipbox{\scriptsize \textbf{Important:} IPFS is a \textbf{distributed file system}, not a transaction tracing tool. It does not analyze blockchain transactions.}
	\end{frame}
	
	% ================================================================
	% SLIDE 13: Screenshot 12 — Data Collection (Q1)
	% ================================================================
	\begin{frame}
		\frametitle{\fontsize{11}{13}\selectfont 12. High-Net-Worth Client Data}
		\begin{block}{\fontsize{8}{10}\selectfont Subject: Data Collection and Preparation}
			\scriptsize
			A compliance officer finds that transactions of an internal high-net-worth client do not match their profile data (collected via KYC). Which data should they review next to assess risk accurately? (Select Two.)
		\end{block}
		\vspace{0.05cm}
		\stepbox{\scriptsize
			\keypoint{Correct Answers:}
			\par Counterparty information.
			\par Transaction velocity metrics.
		}
		\vspace{0.05cm}
		\scriptsize
		\hardconcept{Core Concept:} KYC profile data is a \textbf{static risk indicator}. To assess \textbf{dynamic risk}, you need \textbf{behavioral indicators}. \textbf{Counterparty information} tells you \textit{who} the client is transacting with — is it an offshore shell company? A high-risk jurisdiction? An FIU-flagged entity? \textbf{Transaction velocity} measures the \textit{rate} and \textit{volume} of transactions — a sudden spike is a red flag for layering or structuring. This is consistent with \textbf{FATF Recommendation 10} (ongoing monitoring) and the \textbf{EU AMLD5} requirement to monitor transactions on a risk-sensitive basis.
		\vspace{0.05cm}
		\tipbox{\scriptsize \textbf{Real-World:} Use the \textbf{5W + 1H} approach: Who, What, Where, When, Why, and How. Velocity and counterparty data answer the "Where" and "How often."}
	\end{frame}
	
	% ================================================================
	% SLIDE 14: Screenshot 13 — Data Collection (Q2)
	% ================================================================
	\begin{frame}
		\frametitle{\fontsize{11}{13}\selectfont 13. False Negatives \& Data Quality}
		\begin{block}{\fontsize{8}{10}\selectfont Subject: Data Collection and Preparation}
			\scriptsize
			An internal review finds a high rate of false negatives in the organization's monitoring system due to outdated KYC files and inconsistent country codes. Which is the most reasonable conclusion?
		\end{block}
		\vspace{0.05cm}
		\stepbox{\scriptsize
			\keypoint{Correct Answer:} Data quality controls are insufficient and must be reinforced at the source level.
		}
		\vspace{0.05cm}
		\scriptsize
		\hardconcept{Core Concept:} \textbf{False negatives} are undetected suspicious activities — the most dangerous type of monitoring failure. Outdated KYC means the client risk profile is wrong (e.g., a PEP is not flagged). Inconsistent country codes mean jurisdiction risk is misassigned. The root cause is \textbf{data quality failure at the point of entry}. Under \textbf{BCBS 239} (Principles for Effective Risk Data Aggregation), banks must have effective data quality controls. Under \textbf{FATF Recommendation 18} (Internal Controls), data integrity is a foundational element.
		\vspace{0.05cm}
		\tipbox{\scriptsize \textbf{Key Principle:} \textbf{GIGO — Garbage In, Garbage Out}. Monitoring systems are only as good as the data they ingest. Fix data at the source, not downstream.}
	\end{frame}
	
	% ================================================================
	% SLIDE 15: Screenshot 14 — Global AFC Standards (Q4)
	% ================================================================
	\begin{frame}
		\frametitle{\fontsize{11}{13}\selectfont 14. Corruption Perceptions Index}
		\begin{block}{\fontsize{8}{10}\selectfont Subject: Global AFC Standards and Guidance}
			\scriptsize
			According to Transparency International guidance, the primary purpose of the Corruption Perceptions Index (CPI) is to:
		\end{block}
		\vspace{0.05cm}
		\stepbox{\scriptsize
			\keypoint{Correct Answer:} rank jurisdictions by perceived levels of public sector corruption relative to others.
		}
		\vspace{0.05cm}
		\scriptsize
		\hardconcept{Core Concept:} The \textbf{CPI} is a \textbf{composite index} that uses 13 surveys from 12 independent institutions. It measures \textbf{perceived corruption}, not actual corruption, because corruption is inherently difficult to measure directly. The CPI is used as a \textbf{benchmark} for country risk assessment. It is referenced in \textbf{UN Convention against Corruption} implementation reviews and is used by many banks in their \textbf{country risk scoring} models. A low CPI score is a key indicator for \textbf{enhanced due diligence} on clients from that jurisdiction.
		\vspace{0.05cm}
		\tipbox{\scriptsize \textbf{Real-World Use:} Integrate CPI scores into your \textbf{EWRA (Enterprise-Wide Risk Assessment)}. A country with a CPI score below 50 should trigger enhanced scrutiny.}
	\end{frame}
	
	% ================================================================
	% SLIDE 16: Screenshot 15 — Global AFC Standards (Q5)
	% ================================================================
	\begin{frame}
		\frametitle{\fontsize{11}{13}\selectfont 15. Using CPI for Due Diligence}
		\begin{block}{\fontsize{8}{10}\selectfont Subject: Global AFC Standards and Guidance}
			\scriptsize
			A multinational enterprise is facing challenges in identifying risks related to political integrity and financial secrecy across several countries where it operates. It wants to improve its risk indicators and align its due diligence with authoritative benchmarks. Which of the following would be the most effective use of Transparency International's tools in this situation?
		\end{block}
		\vspace{0.05cm}
		\stepbox{\scriptsize
			\keypoint{Correct Answer:} Referencing the Corruption Perceptions Index (CPI) to assess perceived levels of corruption across jurisdictions.
		}
		\vspace{0.05cm}
		\scriptsize
		\hardconcept{Core Concept:} This is about \textbf{integrating authoritative external benchmarks} into the due diligence process. The CPI provides an \textbf{objective, third-party risk indicator} that is widely recognized by regulators. Using the CPI addresses \textbf{FATF Recommendation 1} (risk-based approach) and aligns with \textbf{OECD Anti-Bribery Convention} requirements. The CPI is also used in \textbf{UNODC} corruption assessments. A comprehensive due diligence program would also use \textbf{Financial Secrecy Index} (FSI) and \textbf{Basel AML Index} as complementary benchmarks.
		\vspace{0.05cm}
		\tipbox{\scriptsize \textbf{Best Practice:} Create a \textbf{country risk matrix} that includes CPI, FSI, FATF mutual evaluation ratings, and the Basel AML Index. Use this matrix to prioritize high-risk jurisdictions.}
	\end{frame}
	
	% ================================================================
	% SLIDE 17: SUMMARY
	% ================================================================
	\begin{frame}
		\frametitle{\fontsize{11}{13}\selectfont SUMMARY: Advanced Concepts Matrix}
		\scriptsize
		\begin{block}{\fontsize{9}{11}\selectfont Regulatory Principles \& Their Application:}
			\vspace{0.05cm}
			\scriptsize
			\begin{tabular}{|p{0.35\textwidth}|p{0.60\textwidth}|}
				\hline
				\keypoint{Concept} & \keypoint{Application} \\
				\hline
				FATF Rec 10 (CDD) & EDD for high-risk clients; UBO identification \\
				\hline
				FATF Rec 15 (Tech) & Centralized monitoring; data quality controls \\
				\hline
				FATF Rec 18 (Internal Controls) & Training programs; role-specific content \\
				\hline
				GDPR Art 5-6 & Balancing AML with data protection; legal basis \\
				\hline
				Three Lines of Defense & First line enforcement; second line oversight \\
				\hline
				Model Risk Mgmt (SR 11-7) & Change management for monitoring systems \\
				\hline
				Blockchain Analytics & Clustering; visualization; privacy challenges \\
				\hline
				Data Quality (BCBS 239) & Source-level controls; GIGO prevention \\
				\hline
				CPI \& Country Risk & Benchmarking corruption for due diligence \\
				\hline
			\end{tabular}
			\vspace{0.05cm}
		\end{block}
		\vspace{0.05cm}
		\tipbox{\scriptsize \textbf{FINAL LESSON:} Every question tests the ability to apply \textbf{regulatory principles} to real-world scenarios. Understand the \textit{why}, and the \textit{what} becomes obvious.}
	\end{frame}
	
	% ================================================================
	% SLIDE 18: THANK YOU
	% ================================================================
	\begin{frame}
		\centering
		\vspace{2.5cm}
		{\fontsize{16}{19}\selectfont \textbf{Thank You}} \\[0.2cm]
		{\fontsize{11}{13}\selectfont Keep learning. Keep growing.} \\[0.15cm]
		\rule{3cm}{0.5pt} \\[0.15cm]
		{\fontsize{8}{10}\selectfont \textit{Compliance is not about perfection — it is about continuous improvement and understanding.}}
		\vspace{0.3cm}
		{\fontsize{7}{9}\selectfont \textcolor{darkgray}{End of Advanced Review}}
	\end{frame}
	
\end{document}